\chapter{KESIMPULAN DAN SARAN}

\section{Kesimpulan}
\label{sec:kesimpulan}

Berdasarkan hasil eksperimen dan analisis yang telah dilakukan, diperoleh
beberapa kesimpulan sebagai berikut.

\begin{enumerate}

    \item
    Penggunaan medium KV-\textit{cache} dapat mengurangi jumlah \textit{token}
    yang digunakan dalam komunikasi antargen dibandingkan medium teks.
    Pada \textit{benchmark} GSM8K, ARC-Challenge, dan HumanEval+, penghematan jumlah
    \textit{token} berada pada rentang $75{,}0\%$--$87{,}8\%$.
    Pengurangan jumlah \textit{token} tersebut juga diikuti oleh waktu
    inferensi yang lebih singkat. Percepatan yang diperoleh berada pada
    rentang sekitar $2{,}5$--$6{,}1$ kali lebih cepat dibandingkan medium teks.
    Pada saat yang sama, akurasi keluaran yang diperoleh dengan medium KV-\textit{cache}
    tidak jauh berbeda dengan medium teks.

    \item
    Di antara metode laten \code{soft}, \code{sample}, dan \code{raw}, \code{moi} menunjukkan hasil
    yang kuat pada gabungan akurasi, efisiensi, dan fidelitas keluaran.
    \code{moi} menghasilkan akurasi rerata sebesar $0{,}850$ dengan 
    jumlah \textit{token} paling sedikit, yaitu $152{,}1$ per soal dan
    waktu inferensi rata-rata $15{,}1$ detik per soal. Dibandingkan medium
    teks yang membutuhkan $974{,}3$ \textit{token} dan $64{,}8$ detik per
    soal, \code{moi} mengurangi jumlah \textit{token} sekitar $84{,}4\%$
    dan waktu inferensi sekitar $76{,}7\%$.

    \item
    Secara keseluruhan, penelitian ini menunjukkan bahwa komunikasi laten
    melalui KV-\textit{cache} dapat menjadi alternatif komunikasi berbasis
    teks yang jauh lebih hemat sumber daya. Namun, hasil yang diperoleh tidak
    hanya ditentukan oleh medium komunikasi, tetapi juga oleh cara informasi
    direpresentasikan dan jenis tugas yang dikerjakan. Hasil geometri
    representasi dan kapasitas kanal laten menunjukkan bahwa keterkaitan
    representasi dengan ruang \textit{embedding} token berkaitan dengan
    kemampuan mempertahankan informasi simbolik, terutama pada tugas yang
    membutuhkan presisi simbolik seperti HumanEval+. Dalam kondisi eksperimen
    penelitian ini, \code{moi} menunjukkan bahwa efisiensi komunikasi dapat
    dicapai sekaligus dengan mempertahankan kinerja yang dekat dengan acuan
    teks.

\end{enumerate}

\section{Saran}
\label{sec:saran}

Berdasarkan hasil dan keterbatasan penelitian ini, beberapa hal yang dapat
dikembangkan pada penelitian selanjutnya adalah sebagai berikut.

\begin{enumerate}

    \item
    Penelitian selanjutnya dapat menggunakan jumlah soal dan jenis
    \textit{benchmark} yang lebih beragam. Penelitian ini menggunakan 100 soal
    pada masing-masing GSM8K, ARC-Challenge, dan HumanEval+. Pengujian pada
    lebih banyak soal dan tugas dengan karakteristik yang berbeda dapat
    membantu melihat apakah pola efisiensi dan akurasi metode laten tetap
    muncul pada kondisi yang lebih beragam.

    \item
    Penelitian selanjutnya dapat mengembangkan metode representasi laten
    selain \code{raw}, \code{soft}, \code{sample}, dan \code{moi}. Pengembangan
    metode representasi lain dapat digunakan untuk mengeksplorasi bentuk
    komunikasi laten yang mampu mempertahankan informasi simbolik dengan
    kebutuhan \textit{token} dan waktu inferensi yang tetap rendah.

    \item
    Penelitian selanjutnya dapat mengembangkan arsitektur bagian keluaran
    model untuk mengurangi ketergantungan terhadap ruang \textit{token} dalam
    proses \textit{decoding} representasi laten. Penggunaan \textit{output
    head}, \textit{adapter}, atau mekanisme pemetaan keluaran lainnya dapat
    dieksplorasi untuk mempertahankan informasi laten sebelum dikonversi
    menjadi \textit{token}.

\end{enumerate}